\documentclass[10pt,a4paper]{amsart}
\usepackage[margin=19mm]{geometry}
\usepackage{amsmath,amssymb,mathtools}
\usepackage[scr=rsfs,cal=euler]{mathalfa}
\usepackage{xfrac}
\usepackage[unicode,hidelinks,bookmarks=false]{hyperref}
\newcommand{\ie}{i.e.\ }
\newcommand{\cf}{cf.\ }
\newcommand{\resp}{respectively\ }
\newcommand{\Subsection}{\subsection}
\providecommand{\nobreakdash}{\nobreak\hbox{-}}
\setlength{\emergencystretch}{2em}
\multlinegap0pt
\usepackage{bm}
\usepackage{bbm}
\usepackage{dsfont}
\usepackage{xspace}
\usepackage{cancel}
\usepackage{mathtools}
\usepackage{amsthm}
\makeatletter
\@ifundefined{equation}{\newtheorem{equation}{Equation}}{}
\@ifundefined{lemma}{\newtheorem{lemma}[equation]{Lemma}}{}
\@ifundefined{theorem}{\newtheorem{theorem}[equation]{Theorem}}{}
\makeatother
\def\B{{\bf B}}
\def\Q{{\bf Q}} \def\Z{{\bf Z}}
\def\R{{\mathrm R}}
\def\lomapr#1{\smash{\mathop{\relbar\joinrel\longrightarrow}\limits^{#1}}}
\newcommand{\FF}{\operatorname{FF} }
\newcommand{\LL}{\operatorname{L} }
\newcommand{\dr}{\operatorname{dR} } 
\newcommand{\hk}{\operatorname{HK} }
\newcommand{\proeet}{\operatorname{pro\acute{e}t} } 
\newcommand{\sd}{{\mathcal{D}}} 
\newcommand{\se}{{\mathcal{E}}}
\title[Duality for $p$-adic geometric pro-\'etale cohomology]{Duality for $p$-adic geometric pro-\'etale cohomology}
\author{Pierre Colmez, Sally Gilles, Wies{\l}awa Nizio{\l}}
\date{Source: arXiv:2411.12163v2 (CC BY 4.0)}
\begin{document}
\maketitle

\noindent\emph{Source and licence.} Duality for $p$-adic geometric pro-\'etale cohomology. Version arXiv:2411.12163v2 (CC BY 4.0), distributed under the Creative Commons Attribution 4.0 International licence, as recorded in the arXiv licence metadata for this exact version and shown on the abstract page https://arxiv.org/abs/2411.12163v2. Full rights notice: licenses/arxiv-2411.12163-CC-BY-4.0.txt.\par\smallskip

\noindent\textit{Editorial scope.} Counted unit: the Lemma whose source label is \texttt{dziecko1} in the article (source0.tex lines 2185--2190, source label \texttt{dziecko1}), delivered together with the article's own complete original proof of it (source0.tex lines 2191--2234, one \texttt{proof} environment of 44 lines). No mathematics is written, repaired or re-derived: statement and proof are the article's own text line for line. Required context retained uncounted: hot1 (lines 1398--1405); kolobrzeg4a (lines 1654--1656). Every definition, display and label of those lines is retained unchanged. Omitted from this package and not redistributed: 1--1397, 1406--1653, 1657--2184, 2235--2791. Nothing inside a retained excerpt is omitted or altered: every retained body is delivered exactly as the article's lines are written, so no omission receipt of any kind accompanies this package. Citations: the retained excerpts carry no citation command at all, so no bibliography entry of the article is transcribed and no reference is redistributed. Reference adaptation: none was needed and none was made; every reference inside the delivered unit resolves inside it, so no unresolved reference or citation is produced. Printed slips kept literal: no printed word, symbol, hypothesis or number of the counted statement or of its proof was altered. Layout and class: the article's own class is \texttt{amsart} with packages including enumitem, amscd, amssymb, xy, fge, todonotes, fontenc, amsfonts, latexsym, url, bbm, hyperref; the wrapper is the lane's pinned \texttt{amsart} editorial wrapper at 10pt on A4, which restates the theorem-like environments the retained text needs and adds the article's own macro definitions that the retained text needs (\texttt{\textbackslash B}, \texttt{\textbackslash FF}, \texttt{\textbackslash LL}, \texttt{\textbackslash Q}, \texttt{\textbackslash R}, \texttt{\textbackslash dr}, \texttt{\textbackslash hk}, \texttt{\textbackslash lomapr}, \texttt{\textbackslash proeet}, \texttt{\textbackslash sd}, \texttt{\textbackslash se}), with extra wrapper packages (bm, bbm, dsfont, xspace, cancel, mathtools). The delivered numbering is the unit's own. Assets: no retained span contains a figure environment, a graphics-inclusion command, a PDF-page inclusion, a table or a TikZ picture; no image file is copied into this package, the package declares no asset dependency and no image file is redistributed with it. Licence: version arXiv:2411.12163v2 of this article is distributed under the Creative Commons Attribution 4.0 International licence, as recorded in the arXiv licence metadata for this exact version and shown on the abstract page https://arxiv.org/abs/2411.12163v2. A full rights notice is delivered at licenses/arxiv-2411.12163-CC-BY-4.0.txt. Characters: every retained excerpt is pure ASCII, so the delivered engine, XeLaTeX with its default Latin Modern text font, reports no missing character.
\begin{theorem}{\rm (Comparison theorem on the $Y_{\FF}$-curve)} \label{hot1} 
Let $X$ be a smooth partially proper variety over $K$.  Let $r\geq 0$. We have  natural, functorial in $S$, and compatible with Frobenius  quasi-isomorphisms in $\sd(\B^{[u,v]}_{S^{\flat},\Box})$ 
and $\sd(\B^{[u,v/p]}_{S^{\flat},\Box})$, respectively:
\begin{align}\label{kin1}
\tau_{\leq r}\R\Gamma_{\proeet,?}(X_S,{\mathbb B}^{[u,v]})(r) & \simeq \tau_{\leq r}[\R\Gamma^{{[u,v]}}_{\hk,?}(X_S,r)\lomapr{\iota_{\hk}} \R\Gamma_{\dr,?}^{{[u,v]}}(X_S,r)],\\
\R\Gamma_{\proeet,?}(X_S,{\mathbb B}^{[u,v/p]})(r) & \simeq \R\Gamma^{{[u,v/p]}}_{\hk,?}(X_S,r).\notag
\end{align}
\end{theorem}

 \begin{lemma}\label{kolobrzeg4a}
The map $ \gamma_{X_S}$ above is a quasi-isomorphism in ${\rm QCoh}(X_{\FF,S^{\flat}})$. 
\end{lemma}

\begin{lemma}\label{dziecko1} 
On $X_{\FF,S^{\flat}}$ we have a natural quasi-isomorphism of solid quasi-coherent sheaves
\begin{align*}
\LL\hskip-.4mm f^*_S\se_{\proeet,?}(X_C,\Q_p) \simeq \se_{\proeet,?}(X_S,\Q_p).
\end{align*}
\end{lemma}

\begin{proof}
It suffices to show that, for a compact rational interval $[u,v]\subset (0,\infty)$, 
the canonical map
$$\R\Gamma_{\proeet,?}(X_C,{\mathbb B}^{[u,v]})\otimes^{\LL_{\Box}}_{\B^{[u,v]}_{C^{\flat}}}
\B_{S^{\flat}}^{ [u,v]}{\to}\R\Gamma_{\proeet,?}(X_S,{\mathbb B}^{[u,v]})$$
is a quasi-isomorphism. Since tensor product commutes with colimits it suffices to do this for  the usual cohomology. 

  By  Theorem \ref{hot1}, it suffices to show the same for the twisted Hyodo-Kato and de Rham cohomologies. Assume thus that $r\geq 2d$
 and let us start  with the Hyodo-Kato cohomology. We want to show that  the  base change map 
(for $\B^{\prime}_S:=\B_{S^{\flat}}^{[u,v]}$)
 $$
 \R\Gamma_{\hk}^{{[u,v]}}(X_C,r)\otimes^{\LL_{\Box}}_{\B^{\prime}_C} \B^{\prime}_{S} \to \R\Gamma_{\hk}^{{[u,v]}}(X_S,r)
  $$
 is a quasi-isomorphism. 
For that,  we may pass to cohomology. Since $H^j_{\hk}(X_C)$ is Fr\'echet (hence flat for the solid tensor product over $\breve{C}$),  this reduces
  to checking that the base change, for $b\geq 0$, 
$$
  (H^b_{\hk}(X_C)\otimes^{\LL_{\Box}}_{\breve{C}}\B^{\prime}_{C,\log})^{N=0}\otimes^{\LL_{\Box}}_{\B^{\prime}_C} \B^{\prime}_S
  \to  (H^b_{\hk}(X_C)\otimes^{\LL_{\Box}}_{\breve{C}}\B^{\prime}_{S,\log})^{N=0}
  $$
  is an isomorphism in $\B^{\prime}_{S, \Box}$. 
  
    Now,
     using the exponential map as in the proof of Lemma \ref{kolobrzeg4a},  we can  reduce  to proving that the base change
  $$
 (H^b_{\hk}(X_C)\otimes^{\LL_{\Box}}_{\breve{C}}\B^{\prime}_C)\otimes^{\LL_{\Box}}_{\B^{\prime}_C} \B^{\prime}_S
  \to  H^b_{\hk}(X_C)\otimes^{\LL_{\Box}}_{\breve{C}}\B^{\prime}_S
  $$
  is an isomorphism in $\B^{\prime}_{S,\Box}$. But this is clear. 
  
  
  
   We pass now to  the de Rham cohomology.  As above it suffices to show that the base change map
   $$
 \R\Gamma_{\dr}^{{[u,v]}}(X_C,r)\otimes^{\LL_{\Box}}_{\B^{\prime}_C} \B^{\prime}_{S} \to \R\Gamma_{\dr}^{{[u,v]}}(X_S,r)
  $$
 is a quasi-isomorphism. But this reduces to showing that the base change maps
 \begin{align*}
& (\Omega^i(X)\otimes^{\Box}_K(\B^+_{\dr}/t^s))
\otimes^{\LL_{\Box}}_{\B^{\prime}_C}\B^{\prime}_S\to  \Omega^i(X)\otimes^{\Box}_K(\B^+_{\dr}(S)/t^s)
%&  H^d_c(X,\Omega^i)\otimes^{\Box}_K(\B^+_{\dr}/t^s)\otimes^{\LL_{\Box}}_{\B^{\prime}_C}\B^{\prime}_S\to H^d_c(X, \Omega^i)\otimes^{\Box}_K(\B^+_{\dr}(S)/t^s)
 \end{align*}
 are isomorphisms. And this is clear.    
\end{proof}

\end{document}
